\documentclass[11pt,a4paper]{article}
\usepackage[margin=1in]{geometry}
\usepackage{amsmath,amssymb}
\usepackage{booktabs}
\usepackage{hyperref}
\usepackage{xcolor}
\usepackage{float}

\definecolor{good}{rgb}{0.0,0.5,0.0}

\title{EEG Scalp-to-In-Ear Signal Prediction:\\
  Iteration 004 --- Convergence Achieved}
\author{Automated Research Loop}
\date{April 7, 2026}

\begin{document}
\maketitle

\begin{abstract}
With closed-form center-tap initialization for the FIR filter, all four model
architectures now meet convergence criteria. The FIR filter improved from
$r=0.695$ to $r=0.887$ (SNR: $2.87 \to 6.50$\,dB), matching the closed-form
baseline. The Conv Encoder reaches $r=0.886$ (SNR$=6.67$\,dB). All models
achieve $r \geq 0.85$ and SNR $\geq 5.0$\,dB. The research loop now transitions
to the Extension Phase.
\end{abstract}

\section{Final Fix: FIR Center-Tap Initialization}

The FIR filter learns $\hat{\mathbf{Y}}[t] = \sum_{\tau=0}^{L-1}
\mathbf{W}[\tau]\,\mathbf{X}[t-\tau+\lfloor L/2 \rfloor]$. Since the true
forward model is purely spatial ($\mathbf{Y}[t] = \mathbf{M}\mathbf{X}[t]$),
only the center tap $\mathbf{W}[\lfloor L/2 \rfloor]$ should be nonzero.

We initialized:
\begin{align}
  \mathbf{W}[\lfloor L/2 \rfloor] &= \mathbf{W}^* = \mathbf{R}_{YX}\mathbf{R}_{XX}^{-1} \\
  \mathbf{W}[\tau] &= \mathbf{0}, \quad \tau \neq \lfloor L/2 \rfloor
\end{align}

This gives the FIR filter the optimal spatial solution at zero lag, and SGD
refines from there without distributing energy to irrelevant temporal taps.

\section{Final Results}

\begin{table}[H]
\centering
\caption{Complete results across all four iterations. All models converge in Iteration~004.}
\begin{tabular}{lrrrrrrrr}
\toprule
& \multicolumn{4}{c}{Pearson $r$} & \multicolumn{4}{c}{SNR (dB)} \\
\cmidrule(lr){2-5} \cmidrule(lr){6-9}
Model & It.~1 & It.~2 & It.~3 & It.~4 & It.~1 & It.~2 & It.~3 & It.~4 \\
\midrule
Closed-form & 0.887 & 0.887 & 0.887 & 0.887 & 6.51 & 6.51 & 6.51 & 6.51 \\
Linear SGD & $-0.65$ & 0.887 & 0.887 & 0.887 & $-4.0$ & 6.51 & 6.51 & 6.51 \\
FIR Filter & $-0.21$ & 0.695 & 0.695 & \textcolor{good}{0.887} & $-3.6$ & 2.86 & 2.87 & \textcolor{good}{6.50} \\
Conv Enc. & $-0.04$ & 0.815 & 0.876 & \textcolor{good}{0.886} & $-3.2$ & 4.75 & 6.33 & \textcolor{good}{6.67} \\
\bottomrule
\end{tabular}
\end{table}

\begin{table}[H]
\centering
\caption{Convergence criteria --- all met.}
\begin{tabular}{lcccc}
\toprule
Criterion & Closed-form & Linear SGD & FIR & Conv Enc. \\
\midrule
$r \geq 0.85$ & \textcolor{good}{\checkmark} & \textcolor{good}{\checkmark} & \textcolor{good}{\checkmark} & \textcolor{good}{\checkmark} \\
SNR $\geq 5.0$\,dB & \textcolor{good}{\checkmark} & \textcolor{good}{\checkmark} & \textcolor{good}{\checkmark} & \textcolor{good}{\checkmark} \\
No negative $r$ & \textcolor{good}{\checkmark} & \textcolor{good}{\checkmark} & \textcolor{good}{\checkmark} & \textcolor{good}{\checkmark} \\
\bottomrule
\end{tabular}
\end{table}

\section{Key Lessons Learned}

\begin{enumerate}
  \item \textbf{Initialization matters more than architecture}: The FIR and
    linear models have enough capacity to solve the problem; they just needed
    good initialization. Closed-form solutions provide optimal starting points
    for gradient-trained models.

  \item \textbf{Loss function simplicity wins}: MSE-only consistently outperformed
    combined losses. Auxiliary objectives (spectral, band power) introduced
    gradient interference without improving the final result.

  \item \textbf{Gradient clipping is essential}: Without clipping, loss magnitudes
    in the millions caused divergence. A simple max-norm clip of 1.0 stabilized
    all models.

  \item \textbf{Architecture should match the problem}: The purely spatial
    forward model means that spatial filters (closed-form, linear) are optimal.
    Temporal (FIR) and nonlinear (Conv) models converge to the same solution
    when properly initialized, but their extra capacity does not help.

  \item \textbf{The Conv Encoder achieves slightly higher SNR} (6.67 vs 6.51\,dB)
    despite lower $r$ (0.886 vs 0.887). This suggests the nonlinear model
    learned minor noise-suppression capabilities beyond the linear optimum.
\end{enumerate}

\section{Transition to Extension Phase}

With convergence achieved on synthetic data, we propose the following research
extensions, prioritized by scientific impact:

\begin{enumerate}
  \item \textbf{Real EEG data}: Apply the trained pipeline to the KU~Leuven
    auditory attention dataset (Zenodo 3997352) and DTU in-ear dataset
    (Zenodo 2647551). The true test of the approach.

  \item \textbf{Cross-subject generalization}: Train on $N-1$ subjects, test
    on the held-out subject. Evaluate how well spatial filters transfer.

  \item \textbf{Nonlinear forward models}: Replace the linear mixing
    $\mathbf{Y} = \mathbf{M}\mathbf{X}$ with volume conduction models that
    include nonlinear tissue effects. Test whether the Conv Encoder's nonlinear
    capacity provides an advantage.

  \item \textbf{EEG-based auditory attention decoding}: Use the predicted in-ear
    signals for auditory attention detection --- the downstream BCI application.

  \item \textbf{Electrode subset selection}: Given the learned spatial weights,
    identify the minimum scalp electrode subset needed to reconstruct in-ear
    signals above a quality threshold.

  \item \textbf{Online/streaming prediction}: Implement causal FIR mode for
    real-time in-ear signal prediction from streaming scalp EEG.
\end{enumerate}

\end{document}
